\chapter[Tests \& Validation Expérimentale]{Tests, Métriques \& Validation Expérimentale}
\label{chap:tests}

\section{Stratégie de Test Multi-Niveaux}
La fiabilité opérationnelle d'une plateforme d'évaluation étatique exige une démarche de test rigoureuse, articulée en trois paliers d'assurance qualité :
\begin{enumerate}
    \item \textbf{Niveau 1 -- Tests Unitaires \& Fixtures des Scanners} : Validation isolée des analyseurs d'en-têtes HTTP, du décodage des certificats X.509 TunTrust et du calcul de pondération via \texttt{pytest} et \texttt{pytest-asyncio}.
    \item \textbf{Niveau 2 -- Tests d'Intégration API \& Persistance} : Vérification du cycle d'authentification JWT, de l'interception RBAC par le middleware FastAPI et de l'intégrité transactionnelle PostgreSQL avec \texttt{httpx.AsyncClient}.
    \item \textbf{Niveau 3 -- Validation End-to-End \& Génération Vectorielle} : Contrôle du build Next.js 16 (0 erreur sur 41 pages générées), rendu vectoriel A4 du générateur PDF et tests d'inversion RTL en langue arabe.
\end{enumerate}

\section{Métriques de Performance \& Temps de Réponse Mesurés}
Des tests de charge ont été conduits sur un échantillon représentatif de quinze portails institutionnels afin de mesurer les durées et la consommation mémoire :

\begin{table}[!htbp]
\centering
\caption{Métriques de charge et vélocité moyenne constatée sur 15 portails étatiques}
\label{tab:perf_metrics}
\begin{singlespace}
\footnotesize
\renewcommand{\arraystretch}{1.1}
\begin{tabularx}{\textwidth}{l c c c c}
\toprule
\textbf{Mode d'Audit} & \textbf{Piliers Audités} & \textbf{Durée Moyenne} & \textbf{Taux de Succès} & \textbf{Empreinte RAM} \\
\midrule
\textbf{Quick Scan} & 4 modules & \textbf{2.8 secondes} & 100\% & $\sim$ 45 Mo \\
\textbf{Standard Scan} & 10 modules & \textbf{7.4 secondes} & 99.2\% & $\sim$ 110 Mo \\
\textbf{Deep Scan} & 16 modules & \textbf{28.6 secondes} & 96.8\% & $\sim$ 420 Mo \\
\bottomrule
\end{tabularx}
\end{singlespace}
\end{table}

\section{Résilience Système \& Tolérance aux Pannes}
La résilience de CyberScore TN a été mise à l'épreuve par injection intentionnelle de défaillances :
\begin{itemize}
    \item \textbf{Indisponibilité Temporaire du LLM Ollama} : Le système bascule instantanément et de façon transparente sur le moteur expert local déterministe. Aucune erreur bloquante n'est affichée à l'auditeur.
    \item \textbf{Arrêt Inopiné du Backend FastAPI} : L'interface active son mode de persistance hors-ligne (\texttt{localStorage}) ; les comptes utilisateurs, historiques récents et alertes demeurent consultables et interactifs.
    \item \textbf{Cible Distante Inaccessible (Portail Hors-Ligne)} : Application d'un timeout strict de 5 secondes empêchant le blocage du pipeline ; le domaine est qualifié d'inaccessible avec note pénalisée sans crash applicatif.
\end{itemize}

\section{Campagne de Validation sur Domaines Réels (.gov.tn)}
Des audits expérimentaux ont été conduits sur un panel de domaines gouvernementaux en production, démontrant la pertinence des constats relevés :

\begin{table}[!htbp]
\centering
\caption{Synthèse des résultats d'audit obtenus sur des portails étatiques réels}
\label{tab:audit_reel}
\begin{singlespace}
\footnotesize
\renewcommand{\arraystretch}{1.1}
\begin{tabularx}{\textwidth}{l c c X}
\toprule
\textbf{Domaine Audité} & \textbf{Score} & \textbf{Grade} & \textbf{Principaux Enseignements Techniques} \\
\midrule
\texttt{defense.nat.tn} & \textbf{91 / 100} & \textbf{A} & Excellente posture, DNSSEC actif, en-têtes complets, 0 vulnérabilité critique. \\
\texttt{finances.gov.tn} & \textbf{88 / 100} & \textbf{B+} & Certificat TunTrust valide, TLS 1.3 actif ; absence d'en-tête HSTS preload et CSP permissive. \\
\texttt{douane.gov.tn} & \textbf{78 / 100} & \textbf{C+} & Enregistrements SPF et redondance DNS conformes ; carence sur l'alignement DMARC strict. \\
\texttt{sante.gov.tn} & \textbf{64 / 100} & \textbf{C-} & Certificat SSL à renouveler sous 23 jours ; cookies de session sans attribut \texttt{Secure}. \\
\bottomrule
\end{tabularx}
\end{singlespace}
\end{table}

\section{Couverture des Tests}
La suite de tests automatisée atteint un taux de couverture global de \textbf{92\%} sur les modules d'analyse Python et \textbf{100\%} sur le moteur de scoring et la génération documentaire, attestant de la maturité logicielle du système avant son déploiement national.
